% Gerado por ferramentas/importar_relatorios.py a partir de MLP/artigo/MLP_EA_bomba.tex.
% Revisado à mão na dissertação: edite este arquivo, não a fonte.
\subsection{Conclusões}
\label{mlp:sec:conclusoes}

Com 56 corridas de um delineamento cruzado de frequência e válvulas, os descritores de \textit{hits} do MOMEA agregados em janelas de 5\,s estimaram a frequência do motor com MAE de 4,2 a 5,6\,Hz para condições e sessões não vistas. MLP, RF e GB tiveram desempenho semelhante. A configuração das válvulas só foi reconhecida quando a própria corrida estava no treino. Com a corrida deixada de fora, a acurácia ficou em 21 a 28\,\% (acaso: 14\,\%). A diferença entre as validações é o principal resultado metodológico: em dados de EA com uma corrida por condição e ordem fixa de ensaio, dividir \textit{hits} ou janelas ao acaso, ou mesmo em blocos dentro da corrida, superestima a capacidade de classificação. Ensaios com repetições aleatorizadas, ganho e limiar fixos e descritores contínuos são necessários para avaliar a detecção de condições ligadas às válvulas e à cavitação.

% =======================================================================================================
